% ============================================================
% ABSTRACT — target: 150--250 words
% ============================================================

Chain-of-thought (CoT) monitoring has emerged as a promising safety mechanism
for understanding the intentions and reasoning processes of large language
models, yet its effectiveness depends on whether models faithfully verbalize
their actual reasoning. Recent work demonstrated that proprietary reasoning
models frequently fail to acknowledge reasoning hints that influenced their
answers---Claude 3.7 Sonnet verbalizes such hints only 25\% of the time, and
DeepSeek-R1 only 39\%---but this evaluation covered just two reasoning models.
We present the first systematic cross-family evaluation of CoT faithfulness in
open-weight reasoning models, testing 11 models spanning 9 architectural
families and ranging from 7B to over 1T parameters. Using 500 multiple-choice
questions from MMLU and GPQA Diamond with six categories of injected reasoning
hints (sycophancy, consistency, visual pattern, metadata, grader hacking, and
unethical information), we measure the rate at which models acknowledge
hint influence in their chain-of-thought when hints successfully alter their
answers. Across 38,500 inference runs, we find that [RESULTS: faithfulness
rates, model-family differences, hint-type patterns, scaling trends]. Our
results reveal [KEY FINDING] and demonstrate that [IMPLICATION for CoT-based
safety monitoring]. These findings establish that CoT faithfulness varies
substantially across the open-weight reasoning model ecosystem and carry direct
implications for the reliability of CoT monitoring in safety-critical
deployments.
